% Глава: 02-linux-access
% Статус: наполнено (первый прогон).
\chapter{Доступ к Linux-окружению}
\label{ch:02-linux-access}

В этой главе мы готовим рабочую среду. Основная модель курса ---
выданные вам виртуальные машины (ВМ) на университетском сервере,
к которым вы подключаетесь по SSH из собственного компьютера.
Локальная установка Linux (WSL, VirtualBox, нативная система)
нужна лишь как запасной вариант и рассмотрена в \ref{sec:local-linux}.

\section{Модель окружения}

Курсу выделяется вычислительный ресурс в виде сервера с хостнеймом
\fpath{pglab.nsu.ru}. На нём разворачиваются виртуальные машины
на базе Ubuntu; каждая ВМ --- отдельная изолированная система,
надёжно отделённая от других (об устройстве виртуализации и контейнеров
см.\ \nameref{ch:19-docker-basics}).

Характеристики выданного окружения:

\begin{itemize}
  \item \textbf{одна ВМ на студента} (по запросу --- больше);
  \item образ Ubuntu \textbf{без дополнительного ПО} --- пакеты ставите сами
    (\nameref{ch:08-packages});
  \item \textbf{20--40 ГБ} дискового пространства;
  \item срок жизни ВМ --- \textbf{до конца семестра};
  \item доступ к ВМ --- только по \textbf{SSH} (паролем или ключом),
    только из внутренней сети; извне --- через \textbf{VPN}
    (заявка на \fpath{support@nsu.ru}).
\end{itemize}

\begin{tipbox}
Работу выполняем в своей виртуальной машине, доступ к которой выдаётся на
сервере \fpath{pglab.nsu.ru}. Подключаемся по SSH:
\cmd{ssh <login>@pglab.nsu.ru} (данные для входа выдаёт семинарист).
Локальное окружение не требуется.
\end{tipbox}

На рис.\ \ref{fig:ssh-access} показана схема подключения.

\begin{figure}[htb]
  \centering
  \begin{tikzpicture}[node distance=1.2cm and 2.4cm, font=\small]
    \node[ext] (laptop) {Ваш компьютер\\[0.15em] Windows / macOS / Linux};
    \node[ext, below=of laptop] (vpn) {VPN-шлюз\\[0.15em] (из внешней сети)};
    \node[process, below=of vpn] (vm) {Ваша ВМ (Ubuntu)\\[0.15em]
        \fpath{pglab.nsu.ru}};
    \node[data, below=of vm] (fs) {Файловая система ВМ\\[0.15em]
        \fpath{/home/<login>}};
    \draw[dottedflow, ->] (laptop.south) -- node[right]{из внешней сети}
        (vpn.north);
    \draw[flow, ->] (vpn.south) -- node[right]{порт 22 (TCP)} (vm.north);
    \draw[flow, <->] (vm) -- (fs);
    \draw[flow, ->] (laptop.east) .. controls +(1.9,0) and +(1.9,0) ..
        node[right, pos=0.5]{SSH-сессия\ (терминал)} (vm.east);
  \end{tikzpicture}
  \caption{Подключение к рабочей ВМ. Из внутренней сети --- напрямую по SSH
    (порт 22); из внешней --- сначала через VPN. Сессия открывает терминал,
    работающий в файловой системе ВМ.}
  \label{fig:ssh-access}
\end{figure}

\section{Подключение по паролю}

SSH --- протокол безопасного удалённого доступа. Клиент
\cmd{ssh} уже установлен практически в любой системе
(проверьте \cmd{ssh -V}).

Проверка, что вы на месте:
\begin{minted}[bgcolor=gray!6]{bash}
ssh <login>@pglab.nsu.ru
\end{minted}

Если семинарист указал нестандартный порт, добавьте \flag{p}:
\begin{minted}[bgcolor=gray!6]{bash}
ssh -p <port> <login>@pglab.nsu.ru
# например: ssh -p 2222 student17@pglab.nsu.ru
\end{minted}

После установки соединения система запросит пароль.
Важная особенность: \textbf{при вводе пароль не отображается вообще}
(ни звёздочками, ни другими символами) --- это нормально.
Просто введите пароль и нажмите Enter.

Завершают сессию командой:
\begin{minted}[bgcolor=gray!6]{bash}
exit        # или нажмите Ctrl+D
\end{minted}

Важно понимать: ваши учебные файлы живут \textbf{на ВМ}, а не на
компьютере. Разрыв SSH-соединения \textbf{не удаляет} сохранённые файлы
--- при следующем входе они на месте (это следствие того, как
файловая система хранилась на диске ВМ, см.\ \nameref{ch:04-filesystem}).
Для входа нужен только доступ к сети.

\section{Подключение по ключу}

Пароль при каждом вводе неудобен и уязвим. Стандартная практика ---
авторизация по \textbf{ключевой паре} (эллиптическая/RSA-криптография).
Механизм асимметричный: у вас есть \emph{приватный ключ} (никому не
показывать) и \emph{публичный ключ} (разрешённый серверу).
При входе сервер шифрует вопрос так, что расшифровать его
может только приватный ключ владельца --- пароль не передаётся.

Генерируем пару на своём компьютере (команда \cmd{ssh-keygen};
разберитесь с ней через \cmd{man ssh-keygen}):
\begin{minted}[bgcolor=gray!6]{bash}
# по умолчанию ключи лежат в ~/.ssh/id_ed25519 и ~/.ssh/id_ed25519.pub
# ed25519: современный и короткий ключ; можно просто принять значения по умолчанию
ssh-keygen -t ed25519
\end{minted}

Копируем публичный ключ на ВМ в файл \fpath{~/.ssh/authorized_keys}.
Если доступ по паролю ещё работает, это делается одной командой:
\begin{minted}[bgcolor=gray!6]{bash}
ssh-copy-id <login>@pglab.nsu.ru
\end{minted}

Если \cmd{ssh-copy-id} недоступен (например, на Windows),
скопируйте публичный ключ \fpath{~/.ssh/id_ed25519.pub}
вручную. Затем допишите его содержимое в файл
\fpath{~/.ssh/authorized_keys} на ВМ
(в \nameref{ch:09-vim} разберём, как править файлы без
графического редактора).

После этого вход без пароля:
\begin{minted}[bgcolor=gray!6]{bash}
ssh <login>@pglab.nsu.ru
\end{minted}

\begin{warnbox}
Приватный ключ --- эквивалент пароля. Не копируйте
\fpath{id_ed25519} (без \fpath{.pub}) в облако и не коммитьте
в Git. Публичный ключ \fpath{id_ed25519.pub} передавать можно.
\end{warnbox}

\subsection{Файл конфигурации SSH}

Постоянно вводить логин, хост и порт не нужно --- опишем их
в конфигурации клиента \fpath{~/.ssh/config} (на вашем компьютере):
\begin{minted}[bgcolor=gray!6]{bash}
Host myvm
    HostName pglab.nsu.ru
    User  <login>
    Port  22
    IdentityFile ~/.ssh/id_ed25519
\end{minted}

Теперь достаточно:
\begin{minted}[bgcolor=gray!6]{bash}
ssh myvm
\end{minted}

\section{Доступ из внешней сети (VPN)}

ВМ доступна только из внутренней сети университета.
Работая из дома или из кафе, сначала подключаемся к VPN
(заявка на \fpath{support@nsu.ru}). После установления VPN-туннеля
подключение \cmd{ssh} работает как обычно --- сервер
становится достижимым по внутренней адресации.

\section{Локальный Linux как запасной вариант}
\label{sec:local-linux}

Если по какой-то причине вы не можете пользоваться сервером,
можно развернуть Linux локально. Основные варианты:
\begin{itemize}
  \item \textbf{WSL} (Windows Subsystem for Linux) --- подсистема,
    позволяющая запускать настоящий Linux из Windows без ВМ;
    устанавливается через «Приложение-программу» или
    \cmd{wsl --install -d Ubuntu} в PowerShell.
  \item \textbf{VirtualBox} --- классическая виртуальная машина;
    загружаем ISO Ubuntu и устанавливаем в ВМ.
  \item \textbf{Нативный Linux} (например, Ubuntu) на компьютере
    или на отдельном разделе.
\end{itemize}

Эти варианты эквивалентны для материала курса: всё, что описано
в остальных главах, выполняется одинаково. Но по умолчанию
мы предполагаем сервер.

\section{Типичные ошибки}

\begin{itemize}
  \item \emph{Connection timed out / No route to host} --- проверка
    сети, адреса и порта; при работе из дома --- установлен ли VPN.
  \item \emph{Permission denied} --- неверный логин или способ входа.
    Не перебирайте пароль многократно (счётчик попыток ведётся);
    покажите сообщение семинаристу.
  \item \emph{Ввёл пароль, но система говорит ``nope''} ---
    убедитесь, что не включён \fpath{CapsLock}, что вы вводите
    пароль для нужного пользователя.
  \item Работа «пропала» после перезагрузки компьютера ---
    вероятно, вы работали в локальной ВМ, а не на сервере;
    на сервере данные сохраняются.
\end{itemize}

\section{Задачи для самопроверки}
\begin{enumerate}
  \item Подключитесь к своей ВМ и выполните
    \cmd{hostname}, \cmd{id}, \cmd{df -h /} --- какие значения
    они выводят и что означает каждый?
  \item Сгенерируйте пару ключей:
    \cmd{ssh-keygen -t ed25519}. Затем настройте \fpath{~/.ssh/config}
    и убедитесь, что \cmd{ssh myvm} работает без запроса пароля.
  \item Объясните своими словами, чем асимметричная авторизация
    по ключу отличается от входа по паролю и почему приватный
    ключ нельзя передавать другим.
  \item Создайте файл \fpath{~/hello.txt} на ВМ, разрыв
    соединение, переподключитесь и убедитесь, что файл на месте.
    Почему он не исчез?
\end{enumerate}
